\chapter{Open Economy General Equilibrium (Ricardian Model)}

\section{Basic Setup}

\begin{itemize}
    \item Two countries: $A, B$
    \item Two goods: $x, y$
    \item One factor: Labor $L$
\end{itemize}

This is a $2 \times 2 \times 1$ Ricardian model.

---

\section{Absolute Advantage}

Unit labor requirements:

\[
\begin{array}{c|cc}
 & A & B \\
\hline
x & a_{Lx}^A = 3 & a_{Lx}^B = 12 \\
y & a_{Ly}^A = 6 & a_{Ly}^B = 4
\end{array}
\]

\begin{itemize}
    \item Country $A$ has absolute advantage in $x$
    \item Country $B$ has absolute advantage in $y$
\end{itemize}

---

\section{Comparative Advantage}

Opportunity cost:

\[
\frac{a_{Lx}^A}{a_{Ly}^A} = \frac{3}{6} = \frac{1}{2}
\]

\[
\frac{a_{Lx}^B}{a_{Ly}^B} = \frac{12}{4} = 3
\]

Thus:
\begin{itemize}
    \item $A$: comparative advantage in $x$
    \item $B$: comparative advantage in $y$
\end{itemize}

---

\section{Production Possibility Curve}

Labor constraint:

\[
a_{Lx} x + a_{Ly} y = L
\]

PPC slope:

\[
\frac{dy}{dx} = -\frac{a_{Lx}}{a_{Ly}}
\]

Thus:
\[
MRT = \frac{a_{Lx}}{a_{Ly}}
\]

---

\section{Autarky Equilibrium}

\subsection{Production}

\[
\max p_x x + p_y y
\quad \text{s.t. } a_{Lx} x + a_{Ly} y = L
\]

FOC:
\[
\frac{p_x}{p_y} = \frac{a_{Lx}}{a_{Ly}}
\]

---

\subsection{Consumption}

\[
\max U(x,y)
\quad \text{s.t. } p_x x + p_y y = I
\]

FOC:
\[
MRS = \frac{p_x}{p_y}
\]

---

\section{Trade Equilibrium}

If:

\[
\left(\frac{p_x}{p_y}\right)^A < \left(\frac{p_x}{p_y}\right)^B
\]

Then world price satisfies:

\[
\left(\frac{p_x}{p_y}\right)^A \leq \left(\frac{p_x}{p_y}\right)^W \leq \left(\frac{p_x}{p_y}\right)^B
\]

---

\section{Specialization}

\begin{itemize}
    \item Country $A$ specializes in $x$
    \item Country $B$ specializes in $y$
\end{itemize}

---

\section{Gains from Trade}

Consumption expands beyond PPC.

\[
C \notin PPC
\]

\[
U^{trade} > U^{autarky}
\]

---

\section{Terms of Trade}

\[
p = \frac{p_x}{p_y}
\]

Trade condition:

\[
p^A < p^W < p^B
\]

---

\section{Trade Balance}

\[
p_x X_P + p_y Y_P = p_x X_C + p_y Y_C
\]

---

\section{Offer Curve}

Offer curve represents excess supply as a function of relative price.

Properties:
\begin{itemize}
    \item Derived from PPC and preferences
    \item Upward sloping
\end{itemize}

---

\section{World Equilibrium}

World equilibrium determined by intersection of offer curves:

\[
EX^A = IM^B
\]

---

\section{Comparative Statics}

\subsection{Price Change}

\[
p \uparrow \Rightarrow X \uparrow
\]

\[
p \downarrow \Rightarrow X \downarrow
\]

---

\section{Multi-Good Extension}

Dornbusch-Fischer-Samuelson (DFS) model:

\[
\frac{a_i^*}{a_i}
\]

Goods ordered by comparative advantage.

---

\section{Wage Determination}

\[
p_i^A = w_A a_i^A
\]

\[
p_i^B = w_B a_i^B
\]

Thus:
\[
\frac{w_A}{w_B} = \frac{a_i^B}{a_i^A}
\]

---

\section{Specialization Condition}

\[
\frac{a_i^*}{a_i} > \frac{w_A}{w_B} \Rightarrow B \text{ produces}
\]

\[
\frac{a_i^*}{a_i} < \frac{w_A}{w_B} \Rightarrow A \text{ produces}
\]

---

\section{Heckscher--Ohlin Preview}

Extension:

\begin{itemize}
    \item Two factors
    \item Factor abundance determines trade
\end{itemize}